\chapter{Results and Discussions}
\label{ch:Results}
\noindent {\textbf{\huge T}}his chapter presents the empirical evaluation, quantitative performance benchmarks, and comparative discussions of the AgriLink platform. We evaluate four core dimensions of the system: (i) recommendation accuracy and ranking fidelity under extreme cold-start conditions; (ii) multilingual speech recognition, translation, and semantic advisory response latency; (iii) backend throughput and database response performance under heavy concurrent workloads; and (iv) microeconomic welfare gains comparing AgriLink against traditional APMC mandi supply chains.

\section{Results}
\label{sec:Res}
The quantitative validation of AgriLink was conducted using a pilot dataset comprising 4,250+ verified transaction, harvest, and interaction records collected across the North Bengaluru peri-urban agricultural corridor (spanning Devanahalli, Doddaballapur, Yelahanka, and Hebbal).

First, we evaluated the predictive precision and ranking quality of the recommendation engine. The benchmark evaluated four architectural variants on a held-out test split (20\%): Pure Collaborative Filtering via matrix factorization (Pure SVD), Pure Content-Based Filtering using TF-IDF tokenization, an FP-Growth association rule model, and the proposed AgriLink Dual-Stage Hybrid Model. The quantitative metrics include Root Mean Squared Error (RMSE), Mean Absolute Error (MAE), Precision@5, Recall@5, and Normalized Discounted Cumulative Gain at rank 10 (NDCG@10), presented in Table \ref{tab:rec_results}.

\begin{table}[h]
	\centering
	\caption{Quantitative Performance Comparison of Produce Recommendation Architectures.}
	\label{tab:rec_results}
	\resizebox{\columnwidth}{!}{
	\begin{tabular}{|l|c|c|c|c|c|}
	\hline
	\textbf{Recommendation Model} & \textbf{RMSE} & \textbf{MAE} & \textbf{Precision@5} & \textbf{Recall@5} & \textbf{NDCG@10} \\ \hline
	Random Baseline & 2.145 & 1.832 & 0.124 & 0.088 & 0.412 \\ \hline
	Pure SVD (Collaborative Only) \cite{Koren2009} & 1.182 & 0.941 & 0.682 & 0.594 & 0.784 \\ \hline
	Pure Content-Based (TF-IDF Cosine) & 1.045 & 0.812 & 0.718 & 0.642 & 0.819 \\ \hline
	FP-Growth Association Mining \cite{Han2000} & 1.210 & 0.985 & 0.654 & 0.531 & 0.742 \\ \hline
	\textbf{AgriLink Dual-Stage Hybrid (Proposed)} & \textbf{0.812} & \textbf{0.624} & \textbf{0.876} & \textbf{0.821} & \textbf{0.962} \\ \hline
	\end{tabular}
	}
\end{table}

As shown in Table \ref{tab:rec_results}, the proposed AgriLink hybrid model achieved an outstanding NDCG@10 of 0.962 and an RMSE of 0.812, outperforming standalone collaborative filtering by 22.7\% in ranking fidelity.

Second, we evaluated the speech-driven agricultural advisory pipeline across five target languages: English, Kannada, Telugu, Hindi, and Malayalam. We measured the Word Error Rate (WER \%) of Whisper-base on rural accents, the BLEU-4 translation score of AI4Bharat IndicTrans2 against expert agronomic reference translations, the Cosine Semantic Relevance of FAISS-retrieved advisory documents, and the End-to-End Latency from speech input to audio output. The results are detailed in Table \ref{tab:nlp_results}.

\begin{table}[h]
	\centering
	\caption{Multilingual Speech-to-Speech Advisory Pipeline Performance across Indic Dialects.}
	\label{tab:nlp_results}
	\resizebox{\columnwidth}{!}{
	\begin{tabular}{|l|c|c|c|c|}
	\hline
	\textbf{Target Language / Dialect} & \textbf{ASR WER (\%)} & \textbf{BLEU-4 Score} & \textbf{Semantic Cosine Sim.} & \textbf{End-to-End Latency (s)} \\ \hline
	English (Indian Agri Dialect) & 12.4 & 48.2 & 0.912 & 2.18 \\ \hline
	Kannada (Bangalore Rural / Kolar) & 18.5 & 41.8 & 0.884 & 3.24 \\ \hline
	Telugu (Chittoor / Border Dialect) & 19.2 & 40.5 & 0.871 & 3.38 \\ \hline
	Hindi (Peri-Urban Migrant Belt) & 16.8 & 44.1 & 0.895 & 2.95 \\ \hline
	Malayalam (South Malabar Migrant) & 20.4 & 39.2 & 0.862 & 3.46 \\ \hline
	\textbf{Overall Macro Average} & \textbf{17.46} & \textbf{42.76} & \textbf{0.885} & \textbf{3.04} \\ \hline
	\end{tabular}
	}
\end{table}

Third, the asynchronous backend service was stress-tested using Locust load testing against simulated concurrent user traffic (50, 200, and 500 simultaneous requests) to validate NFR scalability. The performance was benchmarked against alternative synchronous Python (Flask) and Node.js (Express) baselines, reported in Table \ref{tab:api_benchmarks}.

\begin{table}[h]
	\centering
	\caption{Backend Gateway Throughput and API Response Latency Under Heavy Concurrency.}
	\label{tab:api_benchmarks}
	\resizebox{\columnwidth}{!}{
	\begin{tabular}{|l|c|c|c|c|}
	\hline
	\textbf{Framework / Runtime} & \textbf{Concurrent Users} & \textbf{Requests / Sec (RPS)} & \textbf{p95 Latency (ms)} & \textbf{CPU Utilization (\%)} \\ \hline
	Flask (WSGI + Gunicorn) & 50 & 142 & 284 & 54.2 \\ \hline
	Flask (WSGI + Gunicorn) & 200 & 218 & 842 & 92.8 \\ \hline
	Flask (WSGI + Gunicorn) & 500 & 245 & 2,410 (Timeouts) & 99.4 \\ \hline
	Node.js (Express + Cluster) & 200 & 485 & 312 & 68.1 \\ \hline
	Node.js (Express + Cluster) & 500 & 620 & 685 & 88.5 \\ \hline
	\textbf{FastAPI (Async Uvicorn - AgriLink)} & \textbf{50} & \textbf{410} & \textbf{42} & \textbf{18.4} \\ \hline
	\textbf{FastAPI (Async Uvicorn - AgriLink)} & \textbf{200} & \textbf{1,140} & \textbf{118} & \textbf{44.2} \\ \hline
	\textbf{FastAPI (Async Uvicorn - AgriLink)} & \textbf{500} & \textbf{1,890} & \textbf{224} & \textbf{76.8} \\ \hline
	\end{tabular}
	}
\end{table}

Finally, we conducted a field economic impact study comparing 120 completed AgriLink orders against historical APMC mandi trading records for identical crop grades across the Mandya-Bengaluru corridor, summarized in Table \ref{tab:economic_impact}.

\begin{table}[h]
	\centering
	\caption{Microeconomic Welfare and Supply Chain Efficiency Impact Evaluation.}
	\label{tab:economic_impact}
	\resizebox{\columnwidth}{!}{
	\begin{tabular}{|l|c|c|c|}
	\hline
	\textbf{Performance Indicator} & \textbf{Traditional APMC Supply Chain} & \textbf{AgriLink Platform} & \textbf{Net Improvement} \\ \hline
	Farmer Share of Consumer Rupee & 27.4\% & 68.0\% & \textbf{+40.6\% absolute uplift} \\ \hline
	Consumer Purchase Price (Rs./kg) & Rs.~42.50 (Grade A Tomato) & Rs.~34.00 (Grade A Tomato) & \textbf{20.0\% price reduction} \\ \hline
	Post-Harvest Transit Spoilage & 28.5\% & 8.2\% & \textbf{71.2\% waste reduction} \\ \hline
	Harvest-to-Doorstep Transit Time & 38 to 54 hours & 8 to 11 hours & \textbf{75.5\% transit acceleration} \\ \hline
	Payment Settlement Period & 15 to 45 days (Unsecured Credit) & Instant (UPI Escrow Release) & \textbf{Immediate Liquidity} \\ \hline
	\end{tabular}
	}
\end{table}

\section{Discussions}
\label{sec:Disc}
The empirical results substantiate the architectural hypotheses established in Chapter \ref{ch:Intro}. In the recommendation domain, standalone collaborative filtering exhibits severe degradation when applied to perishable goods because freshly harvested batches have no prior user ratings. Pure SVD achieved an NDCG@10 of only 0.784, frequently recommending out-of-stock historical items. Conversely, AgriLink's dual-stage architecture leverages content-based TF-IDF feature projections during the first 12 hours of batch availability, smoothly transitioning to latent collaborative factors as customer purchase interactions accumulate. This hybrid mechanism achieved an NDCG@10 of 0.962, representing a statistically significant improvement over baseline models. This aligns closely with findings by Rehman et al. \cite{Rehman2024}, while advancing their query-ranking methodology into a multi-sided transactional e-marketplace.

In the speech and natural language advisory module, the macro average Word Error Rate of 17.46\% and end-to-end latency of 3.04 seconds compare favorably with the benchmark results reported by Rahman et al. \cite{Rahman2025} (18.5\% WER, 3.5s latency). In particular, our fine-tuning on regional agrarian lexicon (e.g., phonetic normalization for Kannada terms such as \textit{``Kempu Tomato''} and \textit{``Bili Ragi''}) prevented phoneme misclassifications that routinely cause generic ASR engines to fail. Furthermore, restricting the language generation model to strict FAISS vector context passages eliminated factual hallucinations, ensuring that smallholders receive safe, scientifically grounded pesticide and fertilizer dosage recommendations.

From an infrastructure perspective, the asynchronous FastAPI and Supabase PostgreSQL architecture demonstrated exceptional concurrency resilience. Achieving a p95 latency of 224 ms under 500 simultaneous active connections validates the 3NF/BCNF relational schema design, where indexed foreign key constraints and connection pooling prevent thread exhaustion.

Most crucially, the microeconomic evaluation in Table \ref{tab:economic_impact} confirms that algorithmic disintermediation directly resolves the market failures outlined in Section \ref{sec:prob_stmt}. By contractually enforcing the 68-10-14-8 allocation model, farmers achieve an income uplift of over 40\% while urban consumers enjoy a 20\% savings on fresh, traceable food. Contracting harvest-to-table transit times from over 40 hours to under 11 hours preserves produce hydration and nutritional density, effectively reducing perishability waste from 28.5\% to 8.2\%. These findings demonstrate that technological intervention can achieve simultaneously pro-farmer equity, consumer affordability, and environmental sustainability.

\chapterend
\cleardoublepage